Una vegada més, Einstein tenia raó. Cent anys després d'haver predit l'existència d'ones gravitatòries en la seva teoria general de la relativitat, han estat detectades, i aquesta detecció ha comportat la concessió del Premi Nobel de Física de l'any 2017. Les ones gravitatòries detectades van ser originades per la col·lisió de dos forats negres.

\figura[0.5]{fig1.png}

Igual que les ones gravitatòries, els forats negres també van ser descrits per la teoria general de la relativitat. Les idees bàsiques relatives als forats negres es poden entendre amb les lleis de Newton.

\begin{apartat}{1}
L'any 1783, noranta-sis anys abans del naixement d'Einstein, l'astrònom John Michell (1724-1793) va publicar que un cos esfèric que tingués la mateixa densitat que el Sol i 500 vegades el radi d'aquest tindria una velocitat d'escapament, des de la seva superfície, superior a la velocitat de la llum. Calculeu la massa del cos i aquesta velocitat d'escapament.
\end{apartat}

\begin{apartat}{1}
Calculeu el mòdul de la intensitat del camp gravitatori que el cos de l'apartat anterior crea a la seva pròpia superfície. Quina força (mòdul, direcció i sentit) fa el cos sobre $1\ \mu\mathrm{g}$ situat a la seva superfície?
\end{apartat}

\dades{%
  $G = 6,67\times 10^{-11}\ \mathrm{N\,m^2\,kg^{-2}}$.\\
  Massa del Sol, $M_\mathrm{S} = 1,99\times 10^{30}\ \mathrm{kg}$.\\
  Radi del Sol, $R_\mathrm{S} = 6,96\times 10^{8}\ \mathrm{m}$.}
